% ============================================================
% Methodology appendix/section -- Computational strategy
% Thesis: Frequency Helmholtz Transfer Operator
% Author: F.E. Kiewiet de Jonge
% ============================================================
%
% Suggested packages if this file is used standalone:
%   \usepackage{amsmath, amssymb}
%   \usepackage{booktabs}
%   \usepackage{graphicx}
%
% This file documents how the computational workload was organized across
% local machines, Wave servers, and ORCD cluster resources.

\section{Computational methodology}
\label{sec:methodology_computational}

The numerical experiments in this thesis combine three computationally
different tasks: finite-difference Helmholtz solves, neural-network training,
and Krylov solver evaluation. These tasks stress different parts of the
available hardware. Sparse Helmholtz solves are memory intensive and often
CPU-bound, neural-network training is GPU-bound, and large experimental
sweeps require reliable long-running job management. For this reason the
computational workflow was designed around the strengths of three available
environments: local development, Wave servers, and ORCD cluster resources.

The guiding principle was to use the smallest machine that could answer the
current question reliably. Local runs were used for code development and
small validation tests. Wave servers were used for interactive GPU training,
tmux-managed experiments, and rapid debugging. ORCD resources were used for
larger sweeps, long wall-clock runs, and experiments that required automated
resubmission or multiple independent jobs.


% ============================================================
\subsection{Separation of computational tasks}
\label{subsec:computational_task_separation}
% ============================================================

The pipeline was divided into four stages:
\begin{enumerate}
    \item dataset generation;
    \item neural-network training;
    \item solver evaluation;
    \item spectral and diagnostic analysis.
\end{enumerate}
Each stage was treated differently because its bottleneck is different.

\begin{table}[h]
\centering
\begin{tabular}{llll}
\toprule
Stage & Main bottleneck & Preferred resource & Reason \\
\midrule
Dataset generation & Sparse linear solves, I/O & Wave or ORCD CPU/GPU nodes &
Many independent samples \\
Training & GPU throughput & Wave GPU or ORCD GPU &
UNet training is accelerator-friendly \\
Solver evaluation & Sparse solves and GMRES & Wave or ORCD &
Measures true solver cost \\
Spectral diagnostics & Dense eigensolves in 1D only & Local or Wave CPU &
Small enough for exact checks \\
\bottomrule
\end{tabular}
\caption{Computational stages and their preferred execution environments.}
\label{tab:computational_stages}
\end{table}

This separation avoids a common failure mode in numerical machine learning:
optimizing the neural-network training loop while ignoring the cost of the
linear solver that the model is supposed to improve. The final evaluation is
therefore always solver-native: the relevant metric is not only validation
loss, but the number of Krylov iterations and the residual curve of the actual
Helmholtz solve.


% ============================================================
\subsection{Local development}
\label{subsec:computational_local}
% ============================================================

Local computation was used for small-scale development. The local workflow
served four purposes:
\begin{enumerate}
    \item checking syntax and import errors;
    \item testing plotting scripts on already generated arrays;
    \item validating one-dimensional formulas against analytical references;
    \item inspecting summaries, figures, and logs before launching expensive
    jobs.
\end{enumerate}
Local runs were intentionally kept small. A script that could not complete
quickly on a small problem was not launched at scale until the failure mode
was understood. This reduced wasted cluster time and made the larger runs more
reproducible.


% ============================================================
\subsection{Wave servers for interactive experiments}
\label{subsec:computational_wave}
% ============================================================

Wave servers were used as the main interactive compute environment. They
provided direct access to GPUs for neural-network training and enough memory
for one-dimensional dense spectral analysis and moderate two-dimensional
solver evaluation. Long-running interactive jobs were launched inside
\texttt{tmux} sessions. This made runs robust to SSH disconnects and allowed
training logs to be inspected without interrupting the process.

The typical Wave workflow was:
\begin{enumerate}
    \item start a named \texttt{tmux} session;
    \item launch one clearly scoped experiment, such as one frequency pair or
    one sensitivity setting;
    \item monitor progress through captured logs;
    \item write outputs to a named result directory;
    \item keep exploratory and corrected pipelines in separate folders.
\end{enumerate}
The use of named sessions was important. For example, corrected
one-dimensional flux-PML experiments were launched in sessions named by the
frequency pair and method. This made it possible to return to a run hours
later and immediately identify what it was doing.

GPU allocation on Wave was treated conservatively. Before launching training,
GPU memory and utilization were checked. A single corrected training run was
assigned to one idle GPU rather than competing with other active experiments.
This reduced variance in training time and avoided accidental interference
between jobs.


% ============================================================
\subsection{ORCD resources for large sweeps}
\label{subsec:computational_orcd}
% ============================================================

ORCD cluster resources were used for larger experiments that were too broad
or too long for interactive management. These included model-capacity sweeps,
longer training budgets, residual-loss sweeps, and repeated benchmark
evaluations over multiple frequency pairs. The ORCD workflow was organized
around explicit sweep specifications rather than manual command histories.

The reason is reproducibility. A sweep specification records the variant, the
frequency pair, training budget, loss weights, model capacity, and output
location. This makes it possible to summarize and re-run experiments without
reconstructing decisions from terminal logs. For long jobs, checkpoints such
as \texttt{last.pt} and \texttt{best.pt} were used so that training could
resume safely after wall-time limits.

The ORCD strategy can be summarized as follows:
\begin{table}[h]
\centering
\begin{tabular}{lll}
\toprule
Mechanism & Purpose & Benefit \\
\midrule
Sweep files & Define experiment families & Reproducible comparisons \\
Checkpoints & Save intermediate training state & Survive wall-time limits \\
Automatic summaries & Compare variants systematically & Avoid cherry-picking \\
Benchmark rendering & Generate solver-evaluation jobs & Keep metrics consistent \\
Separate output folders & Preserve exploratory history & Avoid overwriting results \\
\bottomrule
\end{tabular}
\caption{Cluster-level mechanisms used to make large experiments reproducible.}
\label{tab:computational_orcd_mechanisms}
\end{table}


% ============================================================
\subsection{Data and storage strategy}
\label{subsec:computational_storage}
% ============================================================

The main data challenge is size. A two-dimensional complex field on a
$512 \times 512$ grid contains more than a quarter million complex values.
Thousands of such samples are required for training. Storing every auxiliary
quantity naively would quickly make the dataset difficult to move, inspect,
and reproduce.

The storage strategy was therefore selective. Training datasets store the
quantities required to reproduce the supervised learning problem: low- and
high-frequency fields, normalization factors, frequencies, and metadata.
Derived quantities such as residuals, modal coefficients, and plots are
recomputed by analysis scripts when needed. This keeps the primary datasets
compact while preserving the ability to regenerate diagnostics.

Every generated dataset includes metadata describing the grid, PML depth,
frequency pair, source distribution, solver type, and random seed. This is
important because the same visual field can arise from different numerical
conventions. In particular, Green-function data, row-scaled PML data, and
corrected flux-PML data must not be mixed without explicit labels.


% ============================================================
\subsection{Optimization of numerical kernels}
\label{subsec:computational_kernels}
% ============================================================

The sparse Helmholtz matrices are assembled once per frequency and reused
across many right-hand sides whenever possible. Direct sparse factorizations
or preconditioner factorizations are also reused within a batch of solves.
This is essential because repeatedly assembling and factorizing the same
operator would dominate the runtime and obscure the actual cost of the warm
start.

For GMRES evaluation, the comparison is made under a fixed numerical setup:
same matrix, same right-hand sides, same preconditioner, same tolerance, and
same maximum iteration budget. Only the initial guess changes. This makes the
iteration-count comparison attributable to the warm start rather than to
differences in solver configuration.

In the one-dimensional spectral pipeline, dense eigendecompositions are used
only because the matrices are small enough. In two dimensions, complete
eigendecomposition is deliberately avoided. Instead, the two-dimensional
analysis uses scalable diagnostics such as initial residual, interior error,
GMRES iteration count, residual curves, and optionally sparse sampling of
extreme eigenvalues.


% ============================================================
\subsection{Experiment tracking and reproducibility}
\label{subsec:computational_tracking}
% ============================================================

Each experiment is placed in a directory whose name records the main design
choice: frequency pair, training approach, PML convention, or sensitivity
setting. Exploratory results are not overwritten by corrected results. For
example, the corrected one-dimensional flux-PML pipeline is stored separately
from the earlier exploratory one-dimensional results. This preserves the
history of the work while making the final methodology easy to identify.

The reproducibility checklist for a serious result is:
\begin{enumerate}
    \item a script or sweep file that launches the run;
    \item a metadata file describing the numerical setup;
    \item a checkpoint for the trained model;
    \item a summary table of solver metrics;
    \item residual curves and initial-error plots;
    \item a clear distinction between exploratory and thesis-quality figures.
\end{enumerate}
This checklist is especially important because the project combines numerical
PDE choices with machine-learning choices. A result is only interpretable if
both sides of the experiment are recorded.


% ============================================================
\subsection{Computational challenges and safeguards}
\label{subsec:computational_challenges}
% ============================================================

The computational difficulty of this project is not incidental. Helmholtz
systems are indefinite and become harder at higher frequency. PML boundaries
make the matrices complex and non-normal. Neural-network validation loss does
not necessarily predict Krylov convergence. Large two-dimensional datasets
are expensive to regenerate. These challenges create several risks:
\begin{table}[h]
\centering
\begin{tabular}{lll}
\toprule
Risk & Consequence & Safeguard \\
\midrule
Wrong discretization convention & Misleading spectra & Validate in 1D first \\
Field loss improves only visually & No solver benefit & Always benchmark GMRES \\
PML contamination & Inflated initial residual & Test zeroing and PML-trained data \\
Large sweeps become untraceable & Irreproducible conclusions & Use sweep specs and metadata \\
GPU jobs interfere & Unstable runtimes & Check utilization and isolate runs \\
\bottomrule
\end{tabular}
\caption{Main computational risks and safeguards.}
\label{tab:computational_risks}
\end{table}

The main methodological safeguard is to end every important claim with a
solver-level metric. Field loss is useful for training, but it is not the
final objective. The final objective is reduced high-frequency FGMRES cost.
Therefore every serious variant must be evaluated by initial residual,
interior error, iteration count, and residual convergence curve.


% ============================================================
\subsection{Summary}
\label{subsec:computational_summary}
% ============================================================

The computational strategy is deliberately layered. Local runs establish that
the code and formulas are correct. Wave servers provide interactive GPU
training and rapid solver debugging. ORCD resources provide reproducible
large-scale sweeps and benchmarks. The same principle governs all three
levels: use the cheapest reliable computation for the question at hand, and
reserve expensive cluster runs for experiments whose design has already been
validated.

This organization makes the project computationally defensible. It recognizes
that Helmholtz warm-start learning is not only a neural-network problem, but a
coupled numerical linear algebra, PDE discretization, and high-performance
computing problem. The experiments are therefore judged by the metric that
matters for the thesis: whether learned low-frequency information reduces the
cost of solving the high-frequency Helmholtz system.

